\pagebreak
\section{Conclusion}
\label{sec:conclusion}

The demand for high-performance storage in \ac{HPC} environments continues to increase, driven by the demands of the increasing scales of datasets in the age of big data and \ac{AI}. While traditional parallel file systems such as Lustre have been commonplace in \ac{HPC}, Ceph has emerged as a compelling alternative due to its scalability, flexibility, and cost-effectiveness through its unified storage architecture and ability to run on commodity hardware.

Our comprehensive evaluation of Ceph in the context of \ac{HPC} demonstrates that Ceph exhibits significant potential as a high-performance file system for \ac{HPC} environments. Its decentralized storage architecture, enabled by \ac{CRUSH} and \ac{RADOS}, allows for high scalability at the data layer, while its software-defined nature enables granular control over data placement and replication strategies without the need for specialized hardware. Additionally, Ceph's support for the \ac{CephFS} interface, which provides a standard \acs{POSIX}-compliant file system interface while maintaining separation of metadata and data paths, allows it to maintain its data-path performance while providing the semantics of a traditional file system.

However, we have identified two primary limitations of Ceph's architecture. Firstly, the performance of the metadata servers exhibits complex scaling characteristics, with relatively low per-\acs{MDS} performance and inconsistent scaling depending on the directory pinning method and layout of client workloads, possibly requiring hand-tuning for maximal performance. This bottleneck could be a potential limiting factor in its use with metadata-heavy operations, such as modern \ac{AI} workloads, which rely on access to large numbers of relatively small files. Secondly, Ceph's use of the traditional \ac{TCP/IP} stack for its data plane resulted in higher latency and lower overall bandwidth compared to the theoretical limits of the underlying hardware, which we especially observed in small-file write operations. While Ceph allows for the use of an \acs{RDMA} messenger as an experimental feature, its poorly maintained state and lack of official support make it an unreliable option for production environments, with inconsistent performance uplifts compared to \ac{TCP}. However, we did observe the latency benefits at small file sizes with the use of the \ac{RDMA} messenger, which demonstrates the potential for future improvements in this area.

Looking forward, the ongoing development of the Crimson \ac{OSD} and SeaStore backend represents a potentially significant performance leap for future versions of Ceph. Through the redesign of the \ac{OSD} and especially the network layer to better leverage modern flash-based storage and networking hardware with libraries such as \ac{SPDK} and \ac{DPDK}, SeaStore has the potential to significantly improve many of our observed weak points in Ceph's overall design through the minimization of latency and more effective use of the existing hardware. A re-evaluation of Ceph's capabilities after the release of the SeaStore backend would be a valuable direction for future work.

Based on our data, we already consider Ceph a promising candidate for use in \ac{HPC} environments, especially for workloads which are not particularly sensitive to metadata performance, such as large-scale data simulations with checkpointing, or archival storage. It should be noted, however, that our large-scale benchmarks were conducted exclusively with the in-memory \ac{OSD} backend (MemStore), which mostly isolated Ceph's architectural overhead, but does not capture the additional characteristics of physical storage media, such as \acp{SSD} and especially \acp{HDD}. Our results should therefore be considered as an upper bound on Ceph's performance, with real-world deployments likely exhibiting different performance characteristics.

The future release of the SeaStore backend, which among other things targets a major Ceph limitation identified in this thesis---networking---may further broaden its applicability to more metadata-sensitive workloads, such as those found in modern \ac{AI} research. However, there already exist solutions to potentially mitigate some of Ceph's metadata limitations, such as the mounting of overlay file systems on top of large files, or the refactoring of workloads to use a different access pattern when possible, which may make Ceph a viable option for some metadata-sensitive workloads in the short term, so long as the applications can be adapted accordingly.